\subsubsection{\Direction}
\label{ssec:designdims:direction}

Tasks form a dependency graph of dependents, tasks awaiting on others, and dependencies, tasks being awaited. Cancellation starts at a specific node, raising the question: how does cancellation proceed through the graph? We call approaches to this question the Direction dimension.

One approach is \emph{\tdown} cancellation, where cancellation is communicated from dependents to dependencies. This approach describes Rust, which leans on the \code|Drop| trait to implement cascading communication during cancellation. In Smol, when a task is cancelled and later dropped, that drops the handles to all of that task's children, which in turn cancels those tasks, and so on through the task graph. This cancellation approach is feasible in Rust because ownership constrains the shape of the dependency graph to a \abbr{dag}, i.e., nodes have a well-defined parent / child relationship, where parents have an owning reference to the child.

Another approach is \emph{\bup} cancellation, where cancellation is communicated from dependencies to dependents. This approach describes Python, which leans on exceptions to communicate cancellations. When a task is cancelled, both Asyncio and Trio will walk the task graph to find its leaf nodes, and raise a cancellation exception at those tasks. That exception is then raised across await points to async callers, which can choose to either catch the exception or allow it to propagate.
% Asyncio and Trio both offer a \emph{shield} primitive which ensures that cancellation cannot search down certain branches of the task graph, thereby not raising an exception into the leaves of those branches.

The final approach is \emph{\simultaneous} cancellation, where cancellation is communicated at once to all transitive dependencies of the cancelled task. This approach describes Swift, which uses a shared cancellation flag. In Swift, convention holds that async functions should observe the cancellation flag and then raise a cancellation exception. A key difference from \bup cancellation is that the flag ensures that a task is aware of being cancelled, even if the cancellation exception was caught and consumed.

\paragraph{Design rationale.} The \tdown approach is a straightforward means of propagating cancellation without needing a mechanism like exceptions, which is especially useful for exception-less languages such as Rust. As previously mentioned, a \tdown approach requires the dependency graph to be a \abbr{dag}, which has well-defined parent/child relationships. The main limitation of the \tdown approach is that it makes custom handling of cancellations --- such as running cleanup code --- difficult: a cancelled task is never resumed, so it has no opportunity to observe its own cancellation, and cleanup is limited to synchronous destructors (e.g., Rust's \code{Drop} trait).

The \bup and \simultaneous approaches both provide the opportunity for handling cancellations. The major difference is that the \bup approach provides individual async functions more leeway to ignore cancellation, irrespective of the calling context. With the \simultaneous approach, every task in a task tree is guaranteed to know if it is cancelled. As implemented in Swift, this approach requires more complexity --- cancellation is represented both formally through the cancellation flag, and conventionally through a cancellation exception observed by async primitives. In exchange, async functions can avoid the possibility of missed cancellations.

% who has the burden of doing the communication?
% in rust, the dependent has the burden of communication
% in python, it is dependency's burden to propagate up
% in swift, every task is communicated to
